\documentclass[11pt,letterpaper]{article}

% --- Packages ---
\usepackage[margin=1in]{geometry}
\usepackage{graphicx}
\usepackage{booktabs}
\usepackage{longtable}
\usepackage{tabularx}
\usepackage{enumitem}
\usepackage{hyperref}
\usepackage{xcolor}
\usepackage{fancyhdr}
\usepackage{listings}
\usepackage{titlesec}
\usepackage[T1]{fontenc}
\usepackage{lmodern}
\usepackage{parskip}
\usepackage{array}
\usepackage{microtype}

% --- Colours ---
\definecolor{codebg}{gray}{0.95}
\definecolor{codeframe}{gray}{0.75}
\definecolor{linkblue}{HTML}{0066CC}
\definecolor{cautionred}{HTML}{CC0000}
\definecolor{notegray}{HTML}{555555}
\definecolor{guicolor}{HTML}{2E7D32}

% --- Hyperlink styling ---
\hypersetup{
    colorlinks=true,
    linkcolor=linkblue,
    urlcolor=linkblue,
    citecolor=linkblue,
}

% --- Code listing style ---
\lstset{
    backgroundcolor=\color{codebg},
    frame=single,
    rulecolor=\color{codeframe},
    basicstyle=\ttfamily\small,
    breaklines=true,
    breakatwhitespace=false,
    columns=fullflexible,
    keepspaces=true,
    showstringspaces=false,
    xleftmargin=4pt,
    xrightmargin=4pt,
    aboveskip=8pt,
    belowskip=8pt,
}

% --- Header / Footer ---
\pagestyle{fancy}
\fancyhf{}
\fancyhead[L]{\textbf{X-CAVATE Protocol}}
\fancyhead[R]{\thepage}
\renewcommand{\headrulewidth}{0.4pt}

% --- Custom commands ---
\newcommand{\caution}[1]{\par\textcolor{cautionred}{\textbf{CAUTION:} #1}\par}
\newcommand{\note}[1]{\par\textcolor{notegray}{\textbf{Note:} #1}\par}
\newcommand{\critical}[1]{\par\textcolor{cautionred}{\textbf{CRITICAL:} #1}\par}
\newcommand{\code}[1]{\texttt{#1}}
\newcommand{\gui}[1]{\textcolor{guicolor}{\textsf{\textbf{#1}}}}
\newcommand{\btn}[1]{\fbox{\textsf{#1}}}
\newcolumntype{L}[1]{>{\raggedright\arraybackslash}p{#1}}
\newcolumntype{C}[1]{>{\centering\arraybackslash}p{#1}}
\tolerance=400
\emergencystretch=2em

% --- Title ---
\title{
    \vspace{-1cm}
    \textbf{X-CAVATE Protocol}\\[6pt]
    \large G-Code Generation from Vascular Network Geometry\\[4pt]
    \normalsize Stage II of the Vascular Network 3D Bioprinting Workflow
}
\author{}
\date{}

\begin{document}

\maketitle
\thispagestyle{fancy}

\noindent This protocol describes how to use X-CAVATE to convert SimVascular b-spline output into collision-free 3D printer G-code. X-CAVATE can be run in three modes: a \textbf{web GUI} (recommended for new users), a \textbf{command-line interface} (CLI), or via \textbf{Docker}. The steps below cover all three.

\tableofcontents
\newpage

%=============================================================================
\section{Part A: Installation (5 minutes)}
%=============================================================================

\subsection{Option 1: Conda (Recommended)}

\begin{enumerate}
    \item Open a terminal (Mac: Terminal.app; Windows: Anaconda Prompt or PowerShell).

    \item Confirm conda is installed:
    \begin{lstlisting}
conda --version
    \end{lstlisting}
    If not installed, download Anaconda or Miniconda from \url{https://docs.conda.io/en/latest/miniconda.html}.

    \item Navigate to the folder where you want to store X-CAVATE and clone the repository:
    \begin{lstlisting}
git clone https://github.com/skylarscottlab/X-CAVATE.git
cd X-CAVATE
    \end{lstlisting}
    Alternatively, download and unzip the repository from the GitHub page (green ``Code'' button $\to$ ``Download ZIP'').

    \item Create and activate the conda environment (installs Python~3.11, all dependencies, and the X-CAVATE package):
    \begin{lstlisting}
conda env create -f environment.yml
conda activate xcavate-gui
    \end{lstlisting}

    \item Verify the installation:
    \begin{lstlisting}
xcavate --help
    \end{lstlisting}
    You should see the full list of command-line parameters.
\end{enumerate}

To remove the environment later: \code{conda env remove -n xcavate-gui}.

\subsection{Option 2: pip install}

\begin{enumerate}
    \item Open a terminal (Mac: Terminal.app; Windows: Anaconda Prompt or PowerShell).

    \item Confirm Python 3.9+ is installed:
    \begin{lstlisting}
python3 --version
    \end{lstlisting}

    \item Clone and enter the repository:
    \begin{lstlisting}
git clone https://github.com/skylarscottlab/X-CAVATE.git
cd X-CAVATE
    \end{lstlisting}

    \item Install X-CAVATE with GUI support (includes numpy, scipy, pandas, plotly, matplotlib, streamlit):
    \begin{lstlisting}
pip install ".[gui]"
    \end{lstlisting}
    If you only need the CLI (no web interface):
    \begin{lstlisting}
pip install .
    \end{lstlisting}

    \item Verify the installation:
    \begin{lstlisting}
xcavate --help
    \end{lstlisting}
\end{enumerate}

\subsection{Option 3: Docker (no Python required)}

\begin{enumerate}
    \item Install Docker Desktop from \url{https://www.docker.com/products/docker-desktop}.

    \item Open a terminal and navigate to the X-CAVATE folder:
    \begin{lstlisting}
cd X-CAVATE
    \end{lstlisting}

    \item Build the Docker image:
    \begin{lstlisting}
docker build -t xcavate .
    \end{lstlisting}

    \item Run the container:
    \begin{lstlisting}
docker run -p 8501:8501 xcavate
    \end{lstlisting}

    \item Open your browser to \textbf{http://localhost:8501}. The X-CAVATE web interface should appear.
\end{enumerate}


%=============================================================================
\section{Part B: Prepare Input Files (5 minutes)}
%=============================================================================

X-CAVATE requires two input files from your SimVascular workflow.

\subsection{B.1 Network File (b-splines)}

This file contains the vessel centerline coordinates exported from SimVascular. Each vessel has a header line followed by coordinate data rows.

\begin{lstlisting}
Vessel: 0, Number of Points: 100
33.4878, 0.0000, 53.5303, 0.0200
33.4115, -0.5285, 52.7311, 0.0200
...
Vessel: 1, Number of Points: 100
15.0600, -37.3300, 11.6600, 0.0190
...
\end{lstlisting}

\begin{itemize}
    \item \textbf{Columns 1--3} (required): x, y, z coordinates in centimeters.
    \item \textbf{Column 4} (optional): vessel radius, needed if using speed calculation.
    \item \textbf{Column 5} (optional): arterial/venous flag (0 = venous, nonzero = arterial), needed for multimaterial printing.
\end{itemize}

\caution{Coordinates must be in centimeters. X-CAVATE internally converts to millimeters.}
\caution{Limit each vessel to no more than 100 coordinates to avoid long runtimes.}

\subsection{B.2 Inlet/Outlet File}

This file identifies which network locations are inlets (where printing begins) and outlets.

\begin{lstlisting}
inlet
21.97, -15.43, 50.00
11.81, 0.00, 21.34
outlet
36.05, -27.19, 50.00
50.00, 0.00, 26.54
\end{lstlisting}

Each section begins with the keyword \code{inlet} or \code{outlet} on its own line, followed by one or more (x, y, z) coordinate triplets.

\note{Inlet/outlet coordinates are matched to the nearest network point automatically using a KD-tree nearest-neighbor search. They do not need to exactly match a coordinate in the network file. A warning is displayed if the nearest match is more than 5.0\,mm away.}

\subsection{B.3 File Placement}

\begin{itemize}
    \item \textbf{GUI or Docker mode:} Files are uploaded through the web interface. No file placement needed.
    \item \textbf{CLI mode:} Place both files in the \code{inputs/} folder inside the X-CAVATE directory (or specify full paths via \code{--network\_file} and \code{--inletoutlet\_file}).
\end{itemize}


%=============================================================================
\section{Part C: Custom G-Code Templates (5 minutes, optional)}
%=============================================================================

Custom G-code templates adapt X-CAVATE's output to your specific printer hardware and software. Each template file contains a G-code snippet that is automatically inserted at the appropriate location in the final G-code output.

The \code{inputs/custom/} folder contains 12 template files with placeholder text. Replace the placeholder text with your own G-code for each section relevant to your printer.

\subsection{C.1 Header}

\begin{longtable}{@{}p{3cm}p{11cm}@{}}
\toprule
\textbf{File} & \textbf{Instructions} \\
\midrule
\code{header\_code.txt} & Replace the placeholder text with your custom G-code for the header section of the final G-code output. Include any code required for the functioning of your printer, such as establishing connections between external pressure boxes and syringes. \\
\bottomrule
\end{longtable}

\subsection{C.2 Single Material Extrusion}

Skip this section if printing multimaterial; proceed to C.3.

\begin{longtable}{@{}L{5cm}p{9cm}@{}}
\toprule
\textbf{File} & \textbf{Instructions} \\
\midrule
\code{start\_extrusion\_code.txt} & Replace with your G-code for starting extrusion. \\
\code{stop\_extrusion\_code.txt} & Replace with your G-code for stopping extrusion. \\
\bottomrule
\end{longtable}

\subsection{C.3 Multimaterial Extrusion}

Skip this section if printing single material; proceed to C.4.

{\footnotesize
\begin{longtable}{@{}L{5.8cm}p{7.8cm}@{}}
\toprule
\textbf{File} & \textbf{Instructions} \\
\midrule
\code{start\_extrusion\_code\_printhead1.txt} & Start extrusion for Printhead~1 (arterial). \\
\code{start\_extrusion\_code\_printhead2.txt} & Start extrusion for Printhead~2 (venous). \\
\code{stop\_extrusion\_code\_printhead1.txt} & Stop extrusion for Printhead~1 (arterial). \\
\code{stop\_extrusion\_code\_printhead2.txt} & Stop extrusion for Printhead~2 (venous). \\
\code{active\_pressure\_printhead1.txt} & Set ink pressure when PH1 is the \textbf{active} printhead. \\
\code{active\_pressure\_printhead2.txt} & Set ink pressure when PH2 is the \textbf{active} printhead. \\
\code{rest\_pressure\_printhead1.txt} & Set ink pressure when PH1 is \textbf{inactive}. \\
\code{rest\_pressure\_printhead2.txt} & Set ink pressure when PH2 is \textbf{inactive}. \\
\bottomrule
\end{longtable}
}% end footnotesize

\subsection{C.4 Dwell (Optional)}

Incorporating time delays (``dwells'') at the first and last node of each print pass deposits additional ink at each junction, improving the connection via a process similar to ``spot welding.''

\begin{longtable}{@{}p{3cm}p{11cm}@{}}
\toprule
\textbf{File} & \textbf{Instructions} \\
\midrule
\code{dwell\_code.txt} & Replace with your G-code for creating a time delay at the start and end of each print pass (e.g., \code{G4 P80} for an 80\,ms dwell). \\
\bottomrule
\end{longtable}

\subsection{C.5 How to Provide Custom G-Code in Each Mode}

\begin{longtable}{@{}p{2cm}p{12cm}@{}}
\toprule
\textbf{Mode} & \textbf{Instructions} \\
\midrule
\textbf{GUI} & Enable the \gui{Custom G-code} toggle in the sidebar's Advanced section. A tabbed editor (Header, Single Material, Multimaterial, Dwell) appears in the main area. Paste your G-code into the text areas. \\
\textbf{CLI} & Edit the \code{.txt} files in \code{inputs/custom/} directly with a text editor, then run with \code{--custom 1}. \\
\textbf{Docker} & \textit{Option A:} Use the GUI editor in the browser. \textit{Option B:} Mount your local \code{inputs/custom/} folder: \code{docker run -p 8501:8501 -v \$(pwd)/inputs/custom:/app/inputs/custom xcavate}, then enable \gui{Custom G-code} in the GUI. \\
\bottomrule
\end{longtable}


%=============================================================================
\section{Part D: Execute X-CAVATE}
%=============================================================================

\subsection{Mode 1: Web GUI}

\begin{enumerate}
    \item Activate the environment and start the server:
    \begin{lstlisting}
conda activate xcavate-gui
streamlit run xcavate/gui/app.py
    \end{lstlisting}
    A browser window opens at \textbf{http://localhost:8501}.

    \item In the \textbf{sidebar}, upload your two input files and configure all parameters (see Section~\ref{sec:gui-sidebar} for a complete reference).

    \item Click the \btn{Run X-CAVATE} button (blue, full-width, at the bottom of the sidebar). This button is disabled until both input files are uploaded.

    \item A progress bar in the main area tracks the pipeline through 7~steps:
    \begin{enumerate}
        \item Reading network and inlet/outlet files
        \item Interpolating network to nozzle resolution
        \item Building adjacency graph and detecting branchpoints
        \item Generating collision-free print passes (DFS or sweep-line)
        \item Post-processing: subdivision, gap closure, overlap
        \item Processing multimaterial passes (if enabled)
        \item Writing output files and G-code
    \end{enumerate}

    \item When complete, the main area displays:
    \begin{itemize}
        \item \textbf{3D Visualization} --- Interactive Plotly scatter plots of the single-material and (if enabled) multimaterial print passes. Rotate, zoom, and inspect paths directly in the browser.
        \item \textbf{Download buttons} --- Organized in three columns:
        \begin{itemize}
            \item \textbf{G-code} column: one download button per G-code file (e.g., \code{gcode\_SM\_pressure.txt}).
            \item \textbf{Coordinate Files} column: download buttons for all coordinate files (\code{x\_coordinates\_SM.txt}, etc.), speed files, and radii files.
            \item \textbf{Changelog} column: download button for \code{changelog.txt}.
        \end{itemize}
    \end{itemize}

    \item After the run completes, \textbf{Print Instructions} appear in the terminal console showing nozzle positioning, G92 zeroing coordinates, and container centering information.
\end{enumerate}

%-----------------------------------------------------------------------------
\subsubsection{GUI Sidebar Reference}
\label{sec:gui-sidebar}
%-----------------------------------------------------------------------------

The sidebar is organized into the following sections. Every widget, its type, default value, and purpose are listed below.

\paragraph{Input Files}

\begin{longtable}{@{}L{3.5cm} l L{7cm}@{}}
\toprule
\textbf{Widget} & \textbf{Type} & \textbf{Description} \\
\midrule
\gui{Network file (.txt)} & File uploader & Upload the network coordinates file (see Part~B.1). Accepts \code{.txt} files only. \\
\gui{Inlet / Outlet file (.txt)} & File uploader & Upload the inlet/outlet file (see Part~B.2). Accepts \code{.txt} files only. \\
\bottomrule
\end{longtable}

\paragraph{Required Parameters}

\begin{longtable}{@{}L{3.5cm} l l L{5cm}@{}}
\toprule
\textbf{Widget} & \textbf{Type} & \textbf{Default} & \textbf{Description} \\
\midrule
\gui{Nozzle diameter (mm)} & Number & 0.500 & Outer diameter of the print nozzle. Controls interpolation density (nozzle radius = half this value). \\
\gui{Container height (mm)} & Number & 10.0 & Height of the print container. Used for nozzle raise/lower calculations. \\
\gui{Decimal places} & Integer & 3 & Number of decimal places for output coordinate rounding (0--10). \\
\gui{Amount up (mm)} & Number & 10.0 & Distance to raise nozzle above the container between print passes. \\
\gui{Printer type} & Selectbox & Pressure & Dropdown with three options: \textbf{Pressure} (syringe-based), \textbf{Positive Ink} (displacement), \textbf{Aerotech} (6-axis controller). \\
\gui{Multimaterial} & Toggle & Off & Enable dual-ink printing. Requires column~5 (artven flag) in the network file. \\
\bottomrule
\end{longtable}

\paragraph{Optional Parameters --- Tolerancing (expandable)}

\begin{longtable}{@{}L{3.5cm} l l L{5cm}@{}}
\toprule
\textbf{Widget} & \textbf{Type} & \textbf{Default} & \textbf{Description} \\
\midrule
\gui{Enable tolerance} & Toggle & Off & Adds a tolerance offset to vessel dimensions. \\
\gui{Tolerance} & Number & 0.00 & Tolerance amount in mm. Only visible when toggle is on. \\
\bottomrule
\end{longtable}

\paragraph{Optional Parameters --- Print Settings (expandable)}

\begin{longtable}{@{}L{3.5cm} l l L{5cm}@{}}
\toprule
\textbf{Widget} & \textbf{Type} & \textbf{Default} & \textbf{Description} \\
\midrule
\gui{Print speed (mm/s)} & Number & 1.000 & Default print speed for constant-speed mode. \\
\gui{Flow (mm$^3$/s)} & Number & 0.161 & Volumetric flow rate for speed calculation. Used to compute per-node speeds: $v = \text{flow} / (\pi r^2)$. \\
\gui{Jog speed (mm/s)} & Number & 5.00 & Speed for XY travel moves between passes. \\
\gui{Jog speed lift (mm/s)} & Number & 0.25 & Speed for initial Z-lift at end of a pass. \\
\gui{Initial lift (mm)} & Number & 0.50 & Distance for the initial lift before jogging to network top. \\
\gui{Dwell start (s)} & Number & 0.080 & Pause duration at the start of each pass (``spot welding''). \\
\gui{Dwell end (s)} & Number & 0.080 & Pause duration at the end of each pass. \\
\bottomrule
\end{longtable}

\paragraph{Optional Parameters --- Geometry (expandable)}

\begin{longtable}{@{}L{3.5cm} l l L{5cm}@{}}
\toprule
\textbf{Widget} & \textbf{Type} & \textbf{Default} & \textbf{Description} \\
\midrule
\gui{Container X (mm)} & Number & 50.0 & Container x-dimension. Used in print instructions for centering. \\
\gui{Container Y (mm)} & Number & 50.0 & Container y-dimension. Used in print instructions for centering. \\
\gui{Scale factor} & Number & 1.000 & Multiplier applied to all network coordinates after import. \\
\gui{Top padding (mm)} & Number & 0.0 & Extra padding above the highest network point. \\
\bottomrule
\end{longtable}

\paragraph{Optional Parameters --- Downsampling (expandable)}

\begin{longtable}{@{}L{3.5cm} l l L{5cm}@{}}
\toprule
\textbf{Widget} & \textbf{Type} & \textbf{Default} & \textbf{Description} \\
\midrule
\gui{Enable downsampling} & Toggle & Off & Reduces coordinate density after processing by keeping every $N$th node while preserving branchpoints/endpoints. \\
\gui{Downsample factor} & Integer & 1 & Keep every $N$th node. Only visible when toggle is on. \\
\bottomrule
\end{longtable}

\paragraph{Optional Parameters --- Multimaterial (expandable)}

These settings only apply when the \gui{Multimaterial} toggle is on.

\begin{longtable}{@{}L{3.8cm} l l L{4.5cm}@{}}
\toprule
\textbf{Widget} & \textbf{Type} & \textbf{Default} & \textbf{Description} \\
\midrule
\gui{Offset X (mm)} & Number & 103.0 & X-distance between printhead nozzle tips. \\
\gui{Offset Y (mm)} & Number & 0.5 & Y-distance between printhead nozzle tips. \\
\gui{Front nozzle} & Integer & 1 & Which nozzle is in front: 1 = venous in front, 2 = venous behind. \\
\gui{Printhead 1 name} & Text & Aa & Name for the arterial printhead (used in G-code). \\
\gui{Printhead 2 name} & Text & Ab & Name for the venous printhead. \\
\gui{Axis 1 name} & Text & A & Z-axis name for the arterial printhead. \\
\gui{Axis 2 name} & Text & B & Z-axis name for the venous printhead. \\
\gui{Resting pressure (psi)} & Number & 0.0 & Pressure when a printhead is inactive. \\
\gui{Active pressure (psi)} & Number & 5.0 & Pressure when a printhead is actively extruding. \\
\bottomrule
\end{longtable}

\paragraph{Optional Parameters --- Positive Ink Displacement (expandable)}

These settings only apply when \gui{Printer type} is set to \textbf{Positive Ink}.

\begin{longtable}{@{}L{3.8cm} l l L{4.5cm}@{}}
\toprule
\textbf{Widget} & \textbf{Type} & \textbf{Default} & \textbf{Description} \\
\midrule
\gui{Start} & Number & 0.00 & Syringe displacement at pass start (SM). \\
\gui{End} & Number & 0.00 & Syringe displacement at pass end (SM). \\
\gui{Use radii} & Toggle & Off & Compute displacement from vessel radii. \\
\gui{Diameter} & Number & 1.000 & Nozzle inner diameter for displacement calc. \\
\gui{Syringe diameter} & Number & 1.000 & Syringe barrel diameter. \\
\gui{Factor} & Number & 1.000 & Multiplier on computed displacement. \\
\gui{Start (arterial)} & Number & 0.00 & Pass start displacement for arterial (MM). \\
\gui{Start (venous)} & Number & 0.00 & Pass start displacement for venous (MM). \\
\gui{End (arterial)} & Number & 0.00 & Pass end displacement for arterial (MM). \\
\gui{End (venous)} & Number & 0.00 & Pass end displacement for venous (MM). \\
\bottomrule
\end{longtable}

\paragraph{Optional Parameters --- Advanced (expandable)}

\begin{longtable}{@{}L{3.5cm} l l L{5cm}@{}}
\toprule
\textbf{Widget} & \textbf{Type} & \textbf{Default} & \textbf{Description} \\
\midrule
\gui{Custom G-code} & Toggle & Off & Enable custom G-code templates. When on, a tabbed editor appears in the main area with four tabs: \textbf{Header}, \textbf{Single Material}, \textbf{Multimaterial}, \textbf{Dwell}. Paste your G-code snippets into the text areas (see Part~C). \\
\gui{Pathfinding algorithm} & Selectbox & DFS & \textbf{DFS}: modified depth-first search (original algorithm). \textbf{Sweep Line}: bottom-up sweep-line approach (alternative). \\
\gui{Overlap nodes} & Integer & 0 & Number of overlap nodes to add at pass junctions (0 = disabled). Improves print continuity at branchpoints. \\
\gui{Overlap algorithm} & Selectbox & Retrace & Only visible when overlap $> 0$. \textbf{Retrace (original)}: scans all previous passes for shared nodes, retraces backwards, appends to end of current pass. \textbf{Consecutive (fast)}: checks only adjacent pass pairs, prepends to start of next pass. \\
\gui{Generate plots} & Toggle & On & Generate interactive 3D Plotly visualizations. \\
\gui{Speed calculation} & Toggle & Off & Compute per-node print speeds based on vessel radii and flow rate: $v = \text{flow} / (\pi r^2)$. Requires column~4 (radius) in the network file. \\
\bottomrule
\end{longtable}

%-----------------------------------------------------------------------------
\subsubsection{Custom G-Code Template Editor (GUI)}
%-----------------------------------------------------------------------------

When the \gui{Custom G-code} toggle is on, the main area shows a tabbed editor with four tabs:

\begin{enumerate}
    \item \textbf{Header tab} --- A single text area labeled \gui{Header code}. Paste G-code that will appear at the top of the output file (e.g., homing, fan control, connection setup).

    \item \textbf{Single Material tab} --- Two text areas side by side:
    \begin{itemize}
        \item \gui{Start extrusion} (left): G-code to begin extrusion at each pass start.
        \item \gui{Stop extrusion} (right): G-code to stop extrusion at each pass end.
    \end{itemize}

    \item \textbf{Multimaterial tab} --- Eight text areas in two sections:
    \begin{itemize}
        \item \textbf{Extrusion Start / Stop}: four text areas for start/stop extrusion on Printhead~1 and Printhead~2.
        \item \textbf{Pressure Control}: four text areas for active/resting pressure on Printhead~1 and Printhead~2. Each has a help tooltip explaining when it applies.
    \end{itemize}

    \item \textbf{Dwell tab} --- A single text area labeled \gui{Dwell code (applied at start and end of each pass)}. Example: \code{G4 P80} for an 80\,ms pause.
\end{enumerate}

%=============================================================================
\subsection{Mode 2: Command-Line Interface}
%=============================================================================

\begin{enumerate}
    \item Open a terminal, activate the environment, and navigate to the X-CAVATE directory:
    \begin{lstlisting}
conda activate xcavate-gui
cd X-CAVATE
    \end{lstlisting}

    \item (If using custom G-code) Edit the template files in \code{inputs/custom/} (see Part~C), then include \code{--custom 1} in the command.

    \item Run X-CAVATE with your parameters. Example for a single-material pressure print:
    \begin{lstlisting}
xcavate \
  --network_file inputs/VesselNetwork.txt \
  --inletoutlet_file inputs/InletsOutlets.txt \
  --nozzle_diameter 0.65 \
  --container_height 50.0 \
  --num_decimals 5 \
  --amount_up 10.0 \
  --printer_type 0 \
  --multimaterial 0 \
  --tolerance_flag 0 \
  --speed_calc 1 \
  --plots 1 \
  --downsample 0 \
  --custom 1 \
  --flow 0.161
    \end{lstlisting}
    Or equivalently: \code{python -m xcavate [same arguments]}

    \item The terminal outputs progress through each pipeline step, followed by \textbf{Print Instructions} with nozzle positioning, G92 zeroing, container centering, and padding values.

    \item Output files are written to \code{outputs/} (or the directory specified by \code{--output\_dir}).
\end{enumerate}

%-----------------------------------------------------------------------------
\subsubsection{Complete CLI Parameter Reference}
\label{sec:cli-ref}
%-----------------------------------------------------------------------------

\paragraph{Required Parameters}

\begin{longtable}{@{}L{3.5cm} l l L{5cm}@{}}
\toprule
\textbf{Parameter} & \textbf{Type} & \textbf{Example} & \textbf{Description} \\
\midrule
\code{--network\_file} & str & \code{inputs/net.txt} & Path to network coordinates file. \\
\code{--inletoutlet\_file} & str & \code{inputs/io.txt} & Path to inlet/outlet file. \\
\code{--nozzle\_diameter} & float & \code{0.65} & Nozzle outer diameter (mm). \\
\code{--container\_height} & float & \code{50} & Container height (mm). \\
\code{--num\_decimals} & int & \code{5} & Decimal places for rounding. \\
\code{--amount\_up} & float & \code{10} & Z-raise between passes (mm). \\
\code{--printer\_type} & int & \code{0} & 0 = Pressure, 1 = Positive Ink, 2 = Aerotech. \\
\code{--multimaterial} & int & \code{0} & 0 = single ink, 1 = dual ink. \\
\code{--tolerance\_flag} & int & \code{0} & 0 = off, 1 = on. \\
\code{--speed\_calc} & int & \code{1} & 0 = constant speed, 1 = radius-based speed. \\
\code{--plots} & int & \code{1} & 0 = no plots, 1 = generate HTML plots. \\
\code{--downsample} & int & \code{0} & 0 = off, 1 = on. \\
\code{--custom} & int & \code{1} & 0 = no custom G-code, 1 = use templates. \\
\bottomrule
\end{longtable}

\paragraph{Optional --- Geometry}

\begin{longtable}{@{}L{3.3cm} l l L{5.5cm}@{}}
\toprule
\textbf{Parameter} & \textbf{Type} & \textbf{Default} & \textbf{Description} \\
\midrule
\code{--container\_x} & float & 50 & Container x-dimension (mm). \\
\code{--container\_y} & float & 50 & Container y-dimension (mm). \\
\code{--scale\_factor} & float & 1.0 & Multiplier on all coordinates. \\
\code{--top\_padding} & float & 0 & Extra padding above network (mm). \\
\code{--tolerance} & float & 0 & Tolerance amount (mm). \\
\bottomrule
\end{longtable}

\paragraph{Optional --- Print Settings}

\begin{longtable}{@{}L{3.3cm} l l L{5.5cm}@{}}
\toprule
\textbf{Parameter} & \textbf{Type} & \textbf{Default} & \textbf{Description} \\
\midrule
\code{--print\_speed} & float & 1.0 & Constant print speed (mm/s). \\
\code{--flow} & float & 0.161 & Volumetric flow rate (mm$^3$/s). Experimentally determined default. \\
\code{--jog\_speed} & float & 5.0 & XY travel speed (mm/s). \\
\code{--jog\_speed\_lift} & float & 0.25 & Z-lift speed (mm/s). \\
\code{--initial\_lift} & float & 0.5 & Initial lift distance (mm). \\
\code{--jog\_translation} & float & 10 & Speed between nozzles for MM (mm/s). \\
\code{--dwell\_start} & float & 0.08 & Dwell at pass start (seconds). \\
\code{--dwell\_end} & float & 0.08 & Dwell at pass end (seconds). \\
\bottomrule
\end{longtable}

\paragraph{Optional --- Multimaterial}

\begin{longtable}{@{}L{3.5cm} l l L{5cm}@{}}
\toprule
\textbf{Parameter} & \textbf{Type} & \textbf{Default} & \textbf{Description} \\
\midrule
\code{--offset\_x} & float & 103 & Printhead X-offset (mm). \\
\code{--offset\_y} & float & 0.5 & Printhead Y-offset (mm). \\
\code{--front\_nozzle} & int & 1 & 1 = venous in front, 2 = behind. \\
\code{--printhead\_1} & str & Aa & Arterial printhead name. \\
\code{--printhead\_2} & str & Ab & Venous printhead name. \\
\code{--axis\_1} & str & A & Arterial Z-axis name in G-code. \\
\code{--axis\_2} & str & B & Venous Z-axis name in G-code. \\
\code{--resting\_pressure} & float & 0 & Inactive printhead pressure (psi). \\
\code{--active\_pressure} & float & 5 & Active printhead pressure (psi). \\
\bottomrule
\end{longtable}

\paragraph{Optional --- Pathfinding and Post-processing}

\begin{longtable}{@{}L{3.8cm} l l L{4.8cm}@{}}
\toprule
\textbf{Parameter} & \textbf{Type} & \textbf{Default} & \textbf{Description} \\
\midrule
\code{--algorithm} & str & dfs & Pathfinding: \code{dfs} or \code{sweep\_line}. \\
\code{--num\_overlap} & int & 0 & Number of overlap nodes at junctions. \\
\code{--overlap\_algorithm} & str & retrace & \code{retrace}: scans all previous passes (original). \code{consecutive}: checks adjacent pairs (faster). \\
\code{--downsample\_factor} & int & 1 & Keep every $N$th node. \\
\code{--close\_sm} & int & 0 & Use SM gap closure file (0/1). \\
\code{--close\_mm} & int & 0 & Use MM gap closure file (0/1). \\
\code{--output\_dir} & str & outputs & Output directory path. \\
\bottomrule
\end{longtable}

\paragraph{Optional --- Positive Ink Displacement}

{\footnotesize
\begin{longtable}{@{}L{4.8cm} l l L{4cm}@{}}
\toprule
\textbf{Parameter} & \textbf{Type} & \textbf{Default} & \textbf{Description} \\
\midrule
\code{--positiveInk\_start} & float & 0 & Syringe displacement at pass start (SM). \\
\code{--positiveInk\_end} & float & 0 & Syringe displacement at pass end (SM). \\
\code{--positiveInk\_radii} & int & 0 & Use radii for displacement (0/1). \\
\code{--positiveInk\_diam} & float & 1 & Nozzle inner diameter (mm). \\
\code{--positiveInk\_syringe\_diam} & float & 1 & Syringe barrel diameter (mm). \\
\code{--positiveInk\_factor} & float & 1 & Displacement multiplier. \\
\code{--positiveInk\_start\_arterial} & float & 0 & Pass start displacement, arterial (MM). \\
\code{--positiveInk\_start\_venous} & float & 0 & Pass start displacement, venous (MM). \\
\code{--positiveInk\_end\_arterial} & float & 0 & Pass end displacement, arterial (MM). \\
\code{--positiveInk\_end\_venous} & float & 0 & Pass end displacement, venous (MM). \\
\bottomrule
\end{longtable}
}% end footnotesize

%=============================================================================
\subsection{Mode 3: Docker}
%=============================================================================

\begin{enumerate}
    \item Build and start the container (if not already running):
    \begin{lstlisting}
docker build -t xcavate .
docker run -p 8501:8501 -v $(pwd)/outputs:/app/outputs xcavate
    \end{lstlisting}

    \item Open \textbf{http://localhost:8501} in your browser.

    \item Follow the same steps as \textbf{Mode~1 (Web GUI)} above. All functionality is identical.

    \item To use pre-written custom G-code template files, mount the \code{inputs/custom/} folder:
    \begin{lstlisting}
docker run -p 8501:8501 \
  -v $(pwd)/inputs/custom:/app/inputs/custom \
  -v $(pwd)/outputs:/app/outputs \
  xcavate
    \end{lstlisting}

    \item Output files are accessible in the mounted \code{outputs/} directory or via the download buttons.
\end{enumerate}


%=============================================================================
\section{Part E: Inspect Output Files}
%=============================================================================

After execution completes, X-CAVATE produces the following outputs.

\subsection{E.1 G-Code Files}

Located in \code{outputs/gcode/}:

\begin{longtable}{@{}L{5cm}p{9cm}@{}}
\toprule
\textbf{File} & \textbf{Description} \\
\midrule
\code{gcode\_SM\_pressure.txt} & Single-material G-code for Pressure printer. \\
\code{gcode\_SM\_positive\_ink.txt} & Single-material G-code for Positive Ink printer. \\
\code{gcode\_SM\_aerotech.txt} & Single-material G-code for Aerotech printer. \\
\code{gcode\_MM\_*.txt} & Multimaterial equivalents (generated if multimaterial is enabled). \\
\bottomrule
\end{longtable}

The filename suffix matches your selected printer type.

\subsection{E.2 Coordinate Files}

Located in \code{outputs/graph/}:

\begin{longtable}{@{}L{4.5cm}p{9.5cm}@{}}
\toprule
\textbf{File} & \textbf{Description} \\
\midrule
\code{x\_coordinates\_SM.txt} & X-coordinates, one line per node, separated by pass. \\
\code{y\_coordinates\_SM.txt} & Y-coordinates per pass. \\
\code{z\_coordinates\_SM.txt} & Z-coordinates per pass. \\
\code{all\_coordinates\_SM.txt} & Combined x, y, z per pass. \\
\code{radii\_list\_SM.txt} & Vessel radii per node per pass. \\
\code{printspeed\_list\_SM.txt} & Computed print speeds per node (if speed\_calc enabled). \\
\code{*\_MM.txt} & Multimaterial equivalents (if enabled). \\
\code{graph.txt} & Full adjacency graph. \\
\code{special\_nodes.txt} & Branchpoints, endpoints, inlets, outlets with coordinates. \\
\bottomrule
\end{longtable}

\subsection{E.3 Visualizations}

Located in \code{outputs/plots/} (when plots are enabled):

\begin{longtable}{@{}L{5.5cm}p{8.5cm}@{}}
\toprule
\textbf{File} & \textbf{Description} \\
\midrule
\code{network\_original.html} & Pre-interpolation network, one trace per vessel. \\
\code{network\_SM.html} & Single-material print passes (post-processing). \\
\code{network\_downsampled\_SM.html} & Downsampled SM passes (if downsampling enabled). \\
\code{network\_SM\_overlap.html} & SM passes with overlap nodes (if overlap $> 0$). \\
\code{network\_MM.html} & Multimaterial passes colored by arterial/venous (if enabled). \\
\bottomrule
\end{longtable}

Open the \code{.html} files in any web browser to rotate, zoom, and inspect the print paths.

\subsection{E.4 Changelog}

The file \code{outputs/changelog.txt} contains a log of all gap closure operations performed during execution.

\subsection{E.5 Print Instructions (Console Output)}

After G-code generation, X-CAVATE prints operator instructions to the terminal (CLI) or console (GUI). These include:

\begin{itemize}
    \item \textbf{Single material:} Starting nozzle position (\code{X\_start}, \code{Y\_start}), G92 zeroing command, G1 command to move to the first print coordinate, Z positioning instructions.
    \item \textbf{Multimaterial calibration:} Step-by-step procedure for measuring X/Y offsets between two nozzle tips using a calibration point.
    \item \textbf{Multimaterial printing:} Dual-nozzle positioning, dual-axis G92 zeroing, Z-alignment procedure.
    \item \textbf{Padding report:} Left, right, back, and front distances from the network bounding box to the container edges.
\end{itemize}


%=============================================================================
\section{Troubleshooting}
%=============================================================================

\begin{longtable}{@{}L{4cm}p{10cm}@{}}
\toprule
\textbf{Problem} & \textbf{Solution} \\
\midrule
\texttt{Runtime instance already exists} (GUI) & A stale Streamlit process is running. Kill it: \code{pkill -f streamlit}, then restart. \\
Inlet/outlet matching warnings & The inlet/outlet coordinates are approximate. X-CAVATE matches to the nearest network point via KD-tree and logs the distance. Verify the matches are reasonable (warning appears if distance $> 5$\,mm). \\
Pipeline hangs on large networks & Reduce the number of coordinates per vessel to 100 or fewer, or enable downsampling. \\
\texttt{ModuleNotFoundError} & Run \code{pip install ".[gui]"} to install all dependencies. \\
Docker container won't start & Ensure Docker Desktop is running and port 8501 is not in use. \\
\bottomrule
\end{longtable}


%=============================================================================
\section{Quick Reference: Mode Comparison}
%=============================================================================

{\footnotesize
\begin{longtable}{@{}L{3.2cm} L{3.2cm} L{3.5cm} L{3.2cm}@{}}
\toprule
\textbf{Feature} & \textbf{GUI} & \textbf{CLI} & \textbf{Docker} \\
\midrule
Python required & Yes (conda or pip) & Yes (conda or pip) & No \\
Custom G-code via files & No (use GUI editor) & Yes (\code{inputs/\allowbreak custom/}) & Optional (mount or GUI) \\
Custom G-code via editor & Yes (tabbed editor) & No & Yes (tabbed editor) \\
Interactive 3D plots & In-browser (after run) & Open \code{.html} files & In-browser (after run) \\
File download & Download buttons & Files in \code{outputs/} & Buttons + mounted volume \\
Print instructions & Console output & Terminal output & Console output \\
Overlap algorithm & Sidebar selectbox & \code{--overlap\_\allowbreak algorithm} & Sidebar selectbox \\
Pathfinding selection & Sidebar selectbox & \code{--algorithm} & Sidebar selectbox \\
Best for & New users, exploration & Scripting, batch jobs & Reproducibility, sharing \\
\bottomrule
\end{longtable}
}% end footnotesize

\end{document}
